\section{Architecture Decision Records}

Each work stream records its most significant decision. All four are in
status \emph{Proposed} (November 2026).

\subsection{ADR-PAY-01: Hold the money at booking and capture it on confirmation}
\begin{description}
  \item[Context] Booking spans the Reservation Service, Payment Coordination,
    the Ledger and a slow external gateway that cannot join a distributed
    transaction. Dates are held for 10 minutes; a failed booking must never
    leave the guest charged (proposed FR11).
  \item[Decision] \textbf{Authorize} when the payment is requested,
    \textbf{capture} on \texttt{Reservation\allowbreak Confirmed}, and \textbf{void} on
    \texttt{Reservation\allowbreak Expired} or \texttt{Reservation\allowbreak Cancelled} before
    capture.
  \item[Alternatives] Charge immediately and refund on failure; a distributed
    transaction; Payment Coordination confirming the reservation itself.
  \item[Rationale] Releasing a hold is free and immediate, while a refund
    costs fees and days; the gateway has no distributed transactions; reacting
    to events keeps the Reservation Service as sole owner of the booking state.
  \item[Consequences] One extra gateway call per booking; capture must happen
    before the authorization expires; states converge within seconds; every
    gateway operation is idempotent.
\end{description}

\subsection{ADR-RES-01: Prevent double bookings with a database exclusion constraint and expiring holds}
\begin{description}
  \item[Context] A listing must never be booked twice for overlapping dates
    under concurrency ($\approx$130 bookings/s at real scale). Guests need
    time to pay, but abandoned payments must not block the calendar (FR02,
    FR04).
  \item[Decision] Every active period is a calendar entry protected by a
    PostgreSQL exclusion constraint on \texttt{(listing\_id, period)}. Booking
    inserts a \texttt{HOLD} that expires after 10 minutes, released lazily, by
    a scheduled message and by a sweeper, all idempotent.
  \item[Alternatives] \texttt{SELECT ... FOR UPDATE} (relies on every code
    path); a distributed cache lock (can fail mid-transaction); a version
    column (does not cover ranges); pay first (may charge for lost dates).
  \item[Consequences] Double bookings are impossible regardless of instances
    or application bugs, and abandoned holds free themselves. In exchange the
    design depends on range types and \texttt{btree\_gist}, a hot listing
    serializes its bookings, and holds hide dates for up to 10 minutes.
\end{description}

\subsection{ADR-PRICE-01: Answer the multiplier synchronously from a precomputed zone snapshot}
\begin{description}
  \item[Context] FR03 requires a calculated price. The Reservation Service
    needs the multiplier while the host saves a listing ($\approx$500\,ms
    timeout) and Payments uses it as margin. Indicators come from sources with
    very different rhythms: search events (60\,\% of traffic), the calendar
    (daily) and an external events API (monthly).
  \item[Decision] Two indicator services precompute demand and context per
    zone and publish them as events to a local snapshot; each request is one
    read of that snapshot and the active model version, which an administrator
    must approve.
  \item[Alternatives] Compute on request calling the indicator services and
    SerpAPI; save the listing without a multiplier and set it later; a fixed
    margin.
  \item[Rationale] The first chains three services and an external API inside
    500\,ms; the second leaves listings without a rate; the third ignores
    demand and does not satisfy FR03.
  \item[Consequences] Indicators may be a day (demand) or a month (context)
    old, and a new model affects listings only when edited. In return the
    engine is never on the booking or search path, survives indicator outages
    with the last snapshot, and every multiplier can be explained from its
    stored inputs.
\end{description}

\subsection{ADR-ID-01: Deliver notifications asynchronously with at-least-once delivery and idempotent consumers}
\begin{description}
  \item[Context] Notifications must not be duplicated, must handle failures
    per channel, and sending a message, opening a case or changing an account
    must not depend on external email, push or SMS providers.
  \item[Decision] Producers publish through the outbox; Notification applies
    preferences and templates and sends through per-channel queues with
    retries, fallback and a dead-letter queue, using the key
    \texttt{(eventId, userId, channel)} inserted with the outbound task.
  \item[Alternatives] Send inside the request; let each service call the
    providers; assume exactly-once delivery from the broker.
  \item[Rationale] Asynchronous delivery isolates provider outages and allows
    backpressure; exactly-once delivery cannot be assumed, so idempotent
    consumers give an exactly-once \emph{effect}; the outbox removes dual
    writes.
  \item[Consequences] Delivery is eventual; every consumer must be
    idempotent; ordering is per aggregate only; operators need tooling for the
    dead-letter queue.
\end{description}
